\documentclass[10pt]{article}
\usepackage[utf8]{inputenc}
\usepackage{geometry}
\usepackage{graphicx}
\usepackage{hyperref}
\usepackage{xcolor}
\usepackage{booktabs}
\usepackage{fancyhdr}

\geometry{letterpaper, margin=0.75in}

% Header/Footer setup
\pagestyle{fancy}
\setlength{\headheight}{14pt}
\fancyhf{}
\rhead{NOAA EMC/EIB - AI Funding Proposal}
\lhead{December 2, 2025}
\cfoot{\thepage}

\title{\textbf{Scaling AI-Assisted Software Engineering for the\\NOAA Global Workflow}}
\author{Enterprise Infrastructure Branch (EIB)}
\date{}

\begin{document}

\maketitle
\thispagestyle{empty}

\section*{Executive Summary}
The NOAA Environmental Modeling Center (EMC) faces increasing complexity in maintaining and evolving the Global Workflow (GFS, GEFS, GDAS). To address this, we propose scaling our pilot AI-assisted software engineering platform to a "beta" cohort of 10-15 scientific developers. This platform leverages the \textbf{Model Context Protocol (MCP)} and \textbf{Retrieval-Augmented Generation (RAG)} to provide deep, context-aware coding assistance that goes far beyond standard autocomplete. By integrating directly with our specific HPC environments (Hera, Hercules, Orion) and enforcing EE2 compliance standards, this system promises to reduce debugging time, ensure architectural integrity, and accelerate scientific innovation.

\section*{The Solution: A Context-Aware MCP-RAG System}
Standard AI coding tools lack knowledge of NOAA's specific infrastructure, coding standards, and the intricate dependencies of the Global Workflow. Our solution bridges this gap using a specialized architecture:

\begin{itemize}
    \item \textbf{Model Context Protocol (MCP):} A standardized way to connect AI models to our internal systems. It allows the AI to "see" our file systems, read our documentation, and understand our specific workflow structures.
    \item \textbf{Hybrid RAG Engine:} We utilize a dual-database approach for unmatched accuracy:
    \begin{itemize}
        \item \textbf{Vector Database (ChromaDB):} Understands the \textit{semantic meaning} of our documentation and code comments (e.g., "How do I validate environment variables?").
        \item \textbf{Graph Database (Neo4j):} Understands the \textit{structural relationships} of our code (e.g., "Which scripts call \texttt{exglobal\_forecast.py} and what dependencies do they share?").
    \end{itemize}
\end{itemize}

This system allows a developer to ask complex questions like, \textit{"How does the C48\_ATM test case interact with the GSI data assimilation step on Hera?"} and receive an answer grounded in our actual codebase and operational procedures, not just general programming knowledge.

\section*{Value Proposition \& Return on Investment (ROI)}
Investing in this infrastructure yields a high ROI by converting hours of manual debugging and compliance checking into automated, instantaneous tasks. We project the following tangible outcomes:

\begin{itemize}
    \item \textbf{Accelerated Software Development Time (10x Efficiency):} By automating routine lookups, boilerplate generation, and compliance checking, we drastically shorten the development lifecycle. This allows high-cost scientific staff to focus on solving complex modeling problems rather than wrestling with syntax, effectively multiplying our workforce's output without adding headcount.
    \item \textbf{Cost Avoidance via Automated Compliance:} The system proactively enforces NWS/HPC EE2 standards, preventing the accumulation of technical debt. This eliminates the need for costly "compliance cleanup" sprints and reduces the risk of delayed operational implementations.
    \item \textbf{Reduced Operational Risk \& Downtime:} By validating scripts against platform-specific configurations (e.g., WCOSS2 vs. Orion paths) during development, we minimize runtime failures in production, protecting valuable HPC allocations and ensuring timely forecast delivery.
    \item \textbf{Faster Onboarding \& Knowledge Retention:} New staff can query the system to understand legacy components immediately, reducing the training burden on senior engineers and preserving institutional knowledge that might otherwise be lost.
\end{itemize}

\subsection*{Projected Monetary ROI}
The annual cost of the Tier 2 beta environment is approximately \$73,400. Assuming a conservative efficiency gain of just 2 hours per week per developer (automating documentation lookups, debugging, and compliance checks) for a cohort of 20 users:
\begin{itemize}
    \item \textbf{Hours Saved:} 20 users $\times$ 2 hours/week $\times$ 48 weeks = 1,920 hours/year.
    \item \textbf{Equivalent Value:} At a conservative fully-loaded rate of \$150/hr, this equates to \textbf{\$288,000} in annual productivity value.
    \item \textbf{Net ROI:} This represents a \textbf{292\% return on investment}, with the system paying for itself if each developer saves just 30 minutes per week.
\end{itemize}

\section*{Resource Requirements for Beta Rollout (10-15 Users)}
To scale from our current prototype to a robust beta environment for 10-15 users, we require specific infrastructure and licensing. We propose a tiered deployment leveraging the Parallel Works ACTIVATE platform for secure, scalable management.

\subsection*{Infrastructure Needs}
\begin{itemize}
    \item \textbf{Compute Environment:} A single shared AWS instance (e.g., m8i.4xlarge, 16 vCPUs) to host the MCP servers and RAG databases for the entire beta cohort. Individual user clusters are not required.
    \item \textbf{Orchestration:} Parallel Works ACTIVATE platform to manage user sessions, security, and cloud resources.
    \item \textbf{AI Services:}
    \begin{itemize}
        \item \textbf{GitHub Copilot:} Licenses for all beta users (integration point for the IDE) at \$10/month/user.
        \item \textbf{LLM API Keys:} Access to Gemini Pro API keys solely for high-volume document ingestion and knowledge base maintenance. Interactive queries are handled via Copilot, so no per-user API keys are needed.
    \end{itemize}
\end{itemize}

\subsection*{Cost Estimates (Annualized)}
\begin{table}[h]
\centering
\begin{tabular}{@{}lll@{}}
\toprule
\textbf{Component} & \textbf{Tier 1 (Pilot - 1 User)} & \textbf{Tier 2 (Beta - 20 Users)} \\ \midrule
AWS Infrastructure & $\sim$\$6,000 & $\sim$\$6,000 \\
ACTIVATE Platform & \$5,000 (Standard) & \$62,000 (1 Std + 19 Base) \\
GitHub Copilot & \$120 (\$10/mo) & \$2,400 (20 users) \\
LLM API (Ingest Only) & $\sim$\$3,000 & $\sim$\$3,000 \\ \midrule
\textbf{Total Estimated} & \textbf{$\sim$\$14,120} & \textbf{$\sim$\$73,400} \\ \bottomrule
\end{tabular}
\caption{Estimated costs for pilot vs. scaled beta deployment.}
\end{table}

\section*{Conclusion}
The MCP-RAG system represents a transformative step for NOAA's software engineering capability. By embedding our institutional knowledge into an intelligent assistant, we empower our scientists to build better, compliant, and more robust forecasting systems faster. We request approval to proceed with the Tier 2 procurement to enable this capability for our core development team.

\end{document}
